% !TeX program = xelatex
% SPDX-License-Identifier: MIT
% Copyright (c) 2026 Ziyu Peng and 刘毅涵
% Author: https://github.com/pzy54154631-hub?tab=repositories
% 第一作者：Ziyu Peng；第二作者：刘毅涵
% 编译：XeLaTeX -> Biber -> XeLaTeX -> XeLaTeX；或 latexmk -xelatex。
% 所有文献、引文、访谈均为虚构教学示例。
\begin{filecontents*}{latex-07-demo.bib}
@book{sample2024,
  author = {Sample, Alice},
  title = {Local Memory: A Fictional Teaching Example},
  date = {2024},
  location = {Example City},
  publisher = {Fictional Teaching Press}
}
@article{example2025,
  author = {Example, Ben},
  title = {Reading Interviews: A Fictional Teaching Example},
  journaltitle = {Fictional Teaching Journal},
  date = {2025},
  volume = {2},
  number = {1},
  pages = {40--52}
}
\end{filecontents*}
\documentclass[UTF8,a4paper,fontset=fandol]{ctexart}
\usepackage[margin=2.5cm]{geometry}
\usepackage{csquotes,enumitem,xurl}
\usepackage[backend=biber,style=authoryear]{biblatex}
\addbibresource{latex-07-demo.bib}
\usepackage[hidelinks]{hyperref}
\title{把来源和正文一起整理好}
\author{Ziyu Peng \and 刘毅涵}
\date{}
\begin{document}
\maketitle
\tableofcontents
\section{概念与材料}\label{sec:materials}
本文只演示写作与排版。所有作者、文献和访谈均为虚构。
这里的“社区”指日常交往的范围。\footnote{这是一条教学自拟的概念限定。}
\textcite[23--25]{sample2024} 用于演示叙述式引用。
括号式引用用于标记出处\parencite[45]{example2025}。
\begin{quote}
材料的缺席，不能直接说明行动没有发生。
\end{quote}
上面的句子由教程自拟，不是文献原话。
\section{匿名材料整理}
参见第~\ref{sec:materials}~节的材料范围。
以下编号、时间与发言均为虚构。
\begin{description}[leftmargin=2em,font=\normalfont\bfseries]
\item[访问者，00:12] 你怎样描述这处公共空间？
\item[P03，00:18] 我会把它叫作碰面的地方。
\item[转写说明] 方括号用于标识删节和说明，避免与原话混淆。
\end{description}
\printbibliography[heading=bibintoc,title={参考文献（虚构教学示例）}]
\appendix
\section{访谈提纲}\label{app:guide}
这份提纲只演示附录编号。正式写作时，应根据研究问题设计问题。
\section*{致谢}
\phantomsection
\addcontentsline{toc}{section}{致谢}
此处演示无编号标题加入目录。
\end{document}
